\bookchapter{Capitolo 67: Lucidatura e pubblicazione — test, debug e approdo sull'App Store}

\chapterintro{L'app è fatta, e adesso? Questo capitolo risponde alle tre questioni del passaggio da “funziona” a “consegnabile”: come fare debug in modo efficiente per trovare i bug, come scrivere test per prevenire le regressioni e come portare l'app sull'App Store perché il mondo intero possa scaricarla. L'approdo sullo store non è solo una cerimonia: è la linea di confine tra un'opera e un prodotto.}

\booksection{I tre strumenti del debug: print, breakpoint, anteprima}

\begin{center}
\begin{tabularx}{\linewidth}{|X|X|X|}
\hline
\textbf{Strumento} & \textbf{A cosa serve} & \textbf{Come si usa} \tabularnewline
\hline
print + console & Leggere i log di esecuzione & Nel codice print("valore=\textbackslash{}(x)") e l'output si legge nella console in basso a Xcode \tabularnewline
Debug con breakpoint & Far fermare il programma su una riga e vedere i valori delle variabili & Tocca accanto al numero di riga per creare il segnalino blu; all'esecuzione il programma si ferma lì e avanzi passo per passo \tabularnewline
Anteprima SwiftUI + debugger & Vedere in tempo reale l'interfaccia e la gerarchia delle viste & Canvas di \#Preview + l'ispettore View Hierarchy a runtime \tabularnewline
\hline
\end{tabularx}
\end{center}

\begin{center}\small\color{gray}I tre strumenti del debug\end{center}

\begin{tcbinfo}{Consiglio: come leggere un crash?}
Quando Swift crasha, la console stampa “Fatal error” con la causa e Xcode si ferma sulla riga incriminata evidenziandola. I crash passati su dispositivo reale si consultano in Xcode Window Organizer Crashes (serve un account a pagamento). Cerca prima di tutto “la prima frase della causa dell'errore + swiftui”: solo una domanda precisa produce una risposta precisa.

\end{tcbinfo}

\booksection{I test: l'assicurazione del codice}

I test a mano non sono affidabili — correggi una funzione e un'altra si rompe senza farsi notare. I test sono la “pillola dei rimpianti” e il “giubbotto antiproiettile”. Su iOS si usa il framework XCTest, in due varianti:

\begin{center}
\begin{tabularx}{\linewidth}{|X|X|X|X|}
\hline
\textbf{Livello} & \textbf{Cosa verifica} & \textbf{Target} & \textbf{Velocità} \tabularnewline
\hline
Test unitari & Logica pura (calcoli, validazioni) & MyFirstAppTests & In secondi \tabularnewline
Test UI & Comportamento reale dell'interfaccia & MyFirstAppUITests & In minuti \tabularnewline
Test manuale & Le sensazioni e l'esperienza di una persona vera & —— & Serve sempre \tabularnewline
\hline
\end{tabularx}
\end{center}

\begin{center}\small\color{gray}I livelli di test iOS (il progetto modello ha già due Test Target predisposti)\end{center}

\begin{lstlisting}[language=Swift]
import XCTest

// Estrai la logica tipo "stringa importo in numero" in una funzione pura
func parseAmount(_ input: String) -> Double? {
    guard let v = Double(input.trimmingCharacters(in: .whitespaces)),
          v > 0 else { return nil }
    return v
}

final class AmountTests: XCTestCase {
    func testEmptyInputReturnsNil() {
        XCTAssertNil(parseAmount(""))
    }
    func testNegativeReturnsNil() {
        XCTAssertNil(parseAmount("-5"))
    }
    func testNormalValueParses() {
        XCTAssertEqual(parseAmount("12.5"), 12.5)
    }
}
\end{lstlisting}

Eseguire i test: CmdU. Quando è tutto verde puoi modificare il codice con serenità. La validazione degli importi e i calcoli di riepilogo del registro delle spese dovrebbero essere coperti da test unitari.

\booksection{Firma e pacchetto: mettere il “sigillo” all'app}

Ogni app iOS deve essere firmata per poter essere installata — prova che “questo pacchetto viene da te e non è stato manomesso”. I rapporti di gerarchia:

\begin{itemize}
  \item Fase di apprendimento: Apple ID gratuito + firma automatica di Xcode (Personal Team), funziona su dispositivo reale ma va rinnovata ogni 7 giorni;
  \item Fase di distribuzione: un account sviluppatore a pagamento (99 dollari/anno) emette certificati e profili di provisioning, indispensabili per lo store e per l'installazione di lungo periodo.
\end{itemize}

Pacchetto: menu di Xcode Product Archive (creare l'archivio), a build conclusa in Organizer Distribute App scegli App Store Connect carica. L'archiviazione richiede come destinazione Generic iOS Device (non il simulatore): è il passaggio che i principianti dimenticano più spesso.

\textbf{Attenzione alle trappole: certificato perso = aggiornamenti in stallo}\par

Se perdi il certificato di firma e la chiave privata, non potrai più inviare aggiornamenti alla vecchia app: gli utenti potranno installare solo l'applicazione da “nuova”. Fai un backup accurato dei certificati (esportazione .p12 dal portachiavi) e custodisci bene la password. Per fortuna Apple offre la “firma automatica” e la gestione dei certificati nel cloud (in Xcode spunta Automatically manage signing), che evita la maggior parte delle insidie manuali.

\booksection{Il percorso completo per arrivare sull'App Store}

Registrare l'account sviluppatore

Preparare gli asset per lo store

Carica la build

Compila info su privacy e revisione

Revisione Apple (1-7 giorni)

Pubblicato

\begin{center}\small\color{gray}Figura 5: i sei passi per l'approdo sull'App Store\end{center}

\begin{itemize}
  \item Registrazione: developer.apple.com, account personale 99 dollari/anno;
  \item Materiali: icona (1024×1024, senza angoli arrotondati — li taglia Apple), screenshot (per le varie dimensioni di iPhone), descrizione, link all'informativa sulla privacy;
  \item Questionario sulla privacy: dichiara con onestà quali dati raccogli — un registro delle spese senza rete può dichiarare "nessun dato raccolto";
  \item Revisione: alla prima sottomissione l'attesa è spesso di 1-7 giorni; se vieni respinto, correggi punto per punto secondo le "note di revisione" — le motivazioni del rifiuto di Apple sono di solito scritte in modo molto specifico, basta metterle in pratica;
  \item TestFlight: prima della pubblicazione, manda la versione di prova agli amici per il beta testing (fino a 10000 persone), raccogli un giro di feedback e poi sottometti ufficialmente: tutto da professionisti.
\end{itemize}

\booksection{Come fare i materiali dello store: la prima impressione vale i download}

La pagina dello store è la vetrina che l'utente vede mentre “gira per i negozi”. Tre materiali decidono la prima impressione:

\begin{itemize}
  \item Icona: 1024×1024, non improvvisare con uno screenshot. Se non hai competenze di design, genera una forma essenziale con uno strumento di IA: è valida solo se, ridotta sul desktop, resta riconoscibile a colpo d'occhio;
  \item Screenshot: almeno 2 set (per le dimensioni principali di iPhone), che mostrino le scene funzionali più importanti; a ognuno una frase in parole semplici ("Registra una voce in un secondo" vale dieci volte "Efficiente soluzione di contabilità");
  \item Descrizione: la prima frase chiarisca "cos'è + per chi è", senza accumulare "il più professionale, il più potente" — la sincerità funziona meglio dell'esagerazione.
\end{itemize}

\begin{tcbpractice}{Mettiti alla prova: "consegna" una volta il registro delle spese}
\booksection{Segnalazione dei crash e piccola lista per l'ottimizzazione delle prestazioni}

Prima di portare l'app sullo store, queste due cose meritano mezz'ora ciascuna:

\booksubsection{Monitoraggio dei crash}

Se non crasha nei tuoi test, non significa che non crashi sul telefono dell'utente. Apple offre rapporti di crash gratuiti in Xcode Organizer (con l'account a pagamento si vedono i log dei crash su dispositivo reale); la soluzione più comoda è integrare un SDK di terze parti (per esempio Firebase Crashlytics). Pubblicare senza monitoraggio dei crash è come guidare a occhi chiusi.

\booksubsection{I quattro controlli delle prestazioni}

\begin{enumerate}
  \item Per gli elenchi usa LazyVStack / List: non impaginare liste lunghe con un normale VStack — il caricamento pigro è la prima lezione di prestazioni;
  \item Non imporre le immagini originali: usa AsyncImage (immagini remote) e la cache di sistema — un originale da 5MB basta a bloccare i modelli più datati;
  \item Niente lavori pesanti sul thread principale: rete e file grandi passano per async/await;
  \item Misura con Instruments: il profiler incluso in Xcode (Time Profiler, controllo delle perdite di memoria) — prima misura, poi ottimizza, non fidarti delle sensazioni.
\end{enumerate}

(1) Aggiungi 3 test unitari alla logica di parseAmount e simili (verdi con CmdU); (2) con un breakpoint cattura con le tue mani un crash provocato da te e risolvilo; (3) crea un pacchetto con Product Archive e usa TestFlight per invitare gli amici al beta testing; (4) se hai energie in più, registrati come sviluppatore e sottometti davvero il registro delle spese allo store. I primi tre passi si possono completare stasera.

\end{tcbpractice}

\begin{tcbsummary}{Riepilogo del capitolo}
\begin{itemize}
  \item I tre strumenti del debug: log con print, breakpoint, anteprima SwiftUI e ispettore delle viste
  \item XCTest: test unitari (CmdU) + test UI + test manuale, già predisposti nel modello
  \item Firma: firma automatica con l'account gratuito (rinnovo ogni 7 giorni); per distribuire serve l'account a pagamento
  \item Pacchetto: Generic iOS Device + Archive + Distribute; il certificato va assolutamente salvato in backup
  \item I sei passi dello store: registrazione materiali caricamento questionario privacy revisione pubblicazione; prima il beta testing con TestFlight
\end{itemize}

\end{tcbsummary}
